\section{Discussion and open questions}\label{sec:discussion}

The empirical results of Section~\ref{sec:results} show that the
SMC${}^2$ filter and controller, composed in receding-horizon form,
produce a closed-loop policy that improves mean autonomic amplitude by
$+17\%$ over a sensible constant baseline on the FSA-v2 model, with no
fatigue-ceiling violations.

Earlier drafts of these notes flagged three analyses as outstanding:
Fisher information rank identifiability, Lyapunov stability, and a
linear--quadratic--Gaussian (LQG) approximation. The first two have
since been carried out and now form Sections~\ref{sec:fim}
and~\ref{sec:lyap}. The third remains as the principal open piece of
analytical work and is described in
Section~\ref{sec:discussion-lqg} below. Three smaller open questions
are then collected in Section~\ref{sec:discussion-other}.

\subsection{Resolved: Fisher information rank
identifiability}\label{sec:discussion-fim}

The pen-and-paper FIM analysis of Section~\ref{sec:fim} identified
three rank-deficient pairs in the original v2 parametrization,
$(\kappa_B,\varepsilon_A)$, $(\mu_F,\mu_{FF})$, and
$(\tau_F,\lambda_A)$. Each pair was shown to enter the one-day-window
likelihood only through a single linear combination, which is the
combinatorial signature of a rank-one rather than rank-two contribution
to the FIM. The Stage~G1 reparametrization rotates each pair into a
strongly-identified combination and a residual that vanishes at the
operating point of Definition~\ref{def:operating-point}; empirically
this restores full FIM rank and contracts the posterior standard
deviations on the residuals by roughly a factor of six (see
Table~\ref{tab:g1-evidence}).

A useful follow-up, identified in Remark~\ref{rem:fim-where}, is to
add a numerical FIM module under
\texttt{smc2fc/diagnostics/fim/} that recomputes the rank and
condition number directly from the JAX dynamics, so that any future
modification of the model triggers an automatic re-check.

\subsection{Resolved: Lyapunov stability of the
SDE}\label{sec:discussion-lyap}

Section~\ref{sec:lyap} carried out the stability analysis in two
parts. The deterministic skeleton, under a constant control
$\Phi \equiv \Phi_c$, has an explicit $A=0$ equilibrium whose local
stability is determined entirely by the sign of the Stuart--Landau
growth rate $\mu(B^*_0,F^*_0)$
(Proposition~\ref{prop:stab-zero}); the bifurcation locus
\eqref{eq:bif-locus} crosses zero at $\Phi_c \approx 7.5$, above which
overtraining-collapse renders the $A=0$ branch stable. For the full
SDE we exhibited a quadratic Lyapunov function
$V = \tfrac12(aB^2+bF^2+cA^2)$ that satisfies a Foster--Lyapunov drift
inequality (Theorem~\ref{thm:lyap}), and concluded exponential
ergodicity at rate
$\rho < \alpha = \min(2/\tau_B,1/\tau_F^{\mathrm{old}})$
(Theorem~\ref{thm:ergodic}). This delivers strong stability of the
controlled process on the operating domain
$\Xcal = [0,1]\times\R_+^2$ for any bounded control schedule.

Remark~\ref{rem:lyap-limits} noted that the constants in
Theorem~\ref{thm:lyap} are loose: they use the simplest quadratic
Lyapunov function and global bounds on the residual factors. A
sharper rate could be obtained from a piecewise Lyapunov function
that switches between the $\mu < 0$ and $\mu > 0$ regimes identified
in Proposition~\ref{prop:stab-zero}, or from a non-quadratic Lyapunov
function tuned to the geometry of the Stuart--Landau cubic.

\subsection{Outstanding: LQG / Riccati
exact-approximate alternative}\label{sec:discussion-lqg}

Section~\ref{sec:mpc-lqg} sketched a linear--quadratic--Gaussian
controller for FSA-v2 under three explicit assumptions:
\begin{enumerate}
  \item[(A1)] linearisation of the drift around an operating point;
  \item[(A2)] replacement of the state-dependent Jacobi/CIR diffusions
        by constant-coefficient Gaussian noise;
  \item[(A3)] replacement of the asymmetric $-\!\int A\,\dd t$
        objective and the soft fatigue barrier by a quadratic
        surrogate cost.
\end{enumerate}
Under (A1)--(A3), the optimal feedback law is the time-varying linear
policy \eqref{eq:lqg-feedback} with gain matrix $K(t) = R^{-1}
B_{\mathrm{lin}}(t)^\top P(t)$, where $P(t)$ solves the differential
Riccati equation \eqref{eq:riccati}. None of this is yet implemented
in the repository.

We identify three concrete pieces of work that would deliver an LQG
companion to the SMC${}^2$ controller:
\begin{enumerate}
  \item Implement automatic linearisation of the drift around a
        user-supplied operating point, using the existing JAX dynamics
        in \texttt{version\_2/models/fsa\_high\_res/\_dynamics.py};
        expose the time-varying $A_{\mathrm{lin}}(t),B_{\mathrm{lin}}(t)$
        of \eqref{eq:lqg-AB} on the project's fifteen-minute grid.
  \item Implement a backward-integration solver for the differential
        Riccati equation \eqref{eq:riccati}.
  \item Build an \texttt{LQGController} that emits the feedback gain
        $K(t)$ of \eqref{eq:lqg-feedback} and a corresponding
        $\Phi$ schedule.
\end{enumerate}
Two scientific questions then become accessible. \emph{Does the LQG
schedule, evaluated as an open-loop control on the FSA-v2 model, achieve
mean amplitude comparable to the SMC${}^2$ schedule on the $84$-day
benchmark of Section~\ref{sec:results}?} If yes, the non-linear
SMC${}^2$ machinery may be unnecessary in steady-state regimes; if
no, the gap quantifies the value of capturing the cubic
Stuart--Landau non-linearity. \emph{Does using the LQG schedule as a
warm-start for the SMC${}^2$ control of Section~\ref{sec:control}
reduce wall-clock time to solution?} A reasonable hypothesis is that
LQG provides an excellent initial $\thetactrl$ in the linear regime
and only marginal benefit where the non-linearity matters;
experimentally testing this is the natural next step.

\subsection{Other open questions}\label{sec:discussion-other}

A few smaller points deserve mention but are not, in the authors'
view, on the same level of importance as the LQG analysis above.

\begin{itemize}
  \item \emph{Horizon length.} The Stage~E5 result uses a fourteen-day
        macrocycle. The Stage~D open-loop result uses eighty-four days.
        It is not yet known how the closed-loop gain scales with horizon
        length; a horizon sweep is straightforward to set up.
  \item \emph{Posterior shape vs.\ posterior mean.} The controller of
        Section~\ref{sec:mpc-pipeline} uses only the posterior mean
        $\bar\theta_{\mathrm{ctrl}}$. The flat-cost-surface argument of
        Section~\ref{sec:mpc-robustness} explains why this is sometimes
        sufficient, but in regimes where the cost surface is curved the
        full posterior carries useful information. A risk-sensitive
        variant that minimises a quantile of the cost rather than the
        mean would be the natural extension.
  \item \emph{Parameter drift across windows.} The Bures--Wasserstein
        bridge of Section~\ref{sec:smc2} produces a sensible
        Gaussian-fit warm-start, but does not protect against
        \emph{slow} parameter drift in the truth (e.g.\ the athlete's
        time constants change with training age). A more principled
        approach would augment the static parameters with a slow
        random-walk component and let the filter track the drift
        explicitly.
\end{itemize}

\medskip
\noindent
With the FIM and Lyapunov analyses now in place, the model has both
empirical validation through the Stage~D and Stage~E5 results of
Section~\ref{sec:results} and analytical guarantees on identifiability
and stability. The remaining theoretical work concentrates on the
LQG / Riccati alternative.
